Now take for @@M0@@ a \uline{flag variety}
@@M1@@
where @@M2@@ is a simply connected semisimple algebraic group over the algebraically closed field @@M3@@, and @@M4@@ is a Borel subgroup [5]. When @@M5@@, the Bruhat cell decomposition of @@M6@@ [5, Exp. 13] implies that the canonical homomorphism
\sgaStableTarget{sga6.e14.eq.unnamed}@@M7@@
is an isomorphism, as is moreover the canonical homomorphism
\sgaStableTarget{sga6.e14.eq.unnamed.occ2}@@M8@@
where @@M9@@ is the compact space of complex points of @@M10@@, endowed with its usual topology induced by the topology of @@M11@@. On the other hand, L. HODGKIN [12] has proved that @@M12@@ is indeed generated by the classes of invertible sheaves on @@M13@@. The same therefore holds for @@M14@@ itself, and it is moreover not difficult to deduce, by a specialization technique using split semisimple groups over @@M15@@ (SGA 3), that the same result holds over every algebraically closed field @@M16@@. Moreover, the cell decomposition of @@M17@@ (or the bijectivity of @@M18@@ when @@M19@@) readily implies that @@M20@@ is torsion-free, so by 4.2 the homomorphism @@M21@@ is an isomorphism. These facts show that @@M22@@ is an isomorphism if and only if @@M23@@ is generated by its elements of degree @@M24@@, or again, when @@M25@@, if and only if the integral cohomology of @@M26@@ is generated by the elements of degree @@M27@@, i.e. if the classifying space @@M28@@ is torsion-free. It is also known [15, p. 217 th. B (1)@@M29@@(2), and 2.5] that this is so if and only if @@M30@@ is "torsion-free", or equivalently if @@M31@@ is isomorphic to a product of groups @@M32@@ and @@M33@@. Thus it suffices to take for @@M34@@ a group @@M35@@ with @@M36@@, or an exceptional simple group, to obtain an example in which @@M37@@ is not an isomorphism.

\sgaStableTarget{sga6.e14.para.4-6}\textbf{4.6.} Observe that, in this example, there is no doubt that the "right" ring is the torsion-free ring @@M38@@, with the cells of the Bruhat decomposition as a basis, a ring also isomorphic to the integral cohomology of @@M39@@ when @@M40@@; on the other hand, the ring @@M41@@ appears to have torsion of a rather pathological nature. This example therefore confirms the view that the invariant @@M42@@ is of little
